\documentclass[10pt,a4paper]{amsart}
\usepackage[margin=19mm]{geometry}
\usepackage{amsmath,amssymb,mathtools}
\usepackage[scr=rsfs,cal=euler]{mathalfa}
\usepackage{xfrac}
\usepackage[unicode,hidelinks,bookmarks=false]{hyperref}
\newcommand{\ie}{i.e.\ }
\newcommand{\cf}{cf.\ }
\newcommand{\resp}{respectively\ }
\newcommand{\Subsection}{\subsection}
\providecommand{\nobreakdash}{\nobreak\hbox{-}}
\setlength{\emergencystretch}{2em}
\multlinegap0pt
\usepackage{dsfont}
\usepackage{amssymb}
\usepackage{braket}
\usepackage{caption}
\usepackage{enumerate}
\usepackage{natbib}
\usepackage[shortlabels]{enumitem}
\usepackage{multirow}
\usepackage{mathtools}
\usepackage{tcolorbox}
\usepackage{booktabs}
\newtheorem{prop}{Proposition}
\newtheorem{lemma}{Lemma}
\title{\vspace*{-2cm}Riemannian-geometric generalizations of quantum fidelities and Bures-Wasserstein distance}
\author{A. Afham, Chris Ferrie}
\date{Source: arXiv:2410.04937v2 (CC BY 4.0)}
\begin{document}
\maketitle

\noindent\emph{Source and licence.} \vspace*{-2cm}Riemannian-geometric generalizations of quantum fidelities and Bures-Wasserstein distance. Version arXiv:2410.04937v2 (CC BY 4.0), distributed by arXiv under the Creative Commons Attribution 4.0 International licence, http://creativecommons.org/licenses/by/4.0/, as recorded in the arXiv licence metadata for this exact version and shown on the abstract page https://arxiv.org/abs/2410.04937v2. Full licence notice: \texttt{licenses/arxiv-2410.04937-CC-BY-4.0.txt}.\par\smallskip

\noindent\textit{Editorial scope.} Counted unit: the article's prop Thm:GBRAlternativeForm (source0.tex lines 589--595). The article's complete original proof of it is delivered with it (source0.tex lines 634--665, one \texttt{proof} environment of 32 lines, which the article places away from the statement). No mathematics is written, repaired or re-derived: the statement and the proof are the article's own lines, in the article's own notation, and no retained line is substituted or re-flowed.  Retained uncounted as the source-local context the proof needs: lines 1858--1868.  Nothing was omitted from any retained span, so the package carries no omission record.  No retained line contains a citation, so the wrapper carries no bibliography and no bibliography entry.  No retained line contains a figure environment, an image-inclusion command or drawing code, so no image is delivered and the package declares no asset dependency.  Editorial formatting only, in addition: an AMS \texttt{amsart} preamble that restates the article's own declarations and macros used by the retained lines, namely the environments \texttt{prop}, \texttt{lemma} on separate counters, together with the standard packages those lines need.  The retained environments print in this document as Proposition 1, Lemma 1 in order of appearance, whereas the article prints the same items with the numbering of its own full document.  The complete native source of the article as distributed in the arXiv e-print for this exact version is bundled unchanged as \texttt{sources/166411/source0.tex}.
    \begin{lemma}\label{Lem:PolarFactorLemma}
Let $P, Q \in \mathbb{P}_d$ and let $U := \operatorname{Pol}\left(P^{\frac 12} Q^{\frac 12}\right)$. Then,
\begin{align}
    U^* P^{\frac 12} Q^{\frac 12} &= \sqrt{Q^{\frac 12} P Q^{\frac 12}} = Q^{\frac 12} P^{\frac 12} U, \label{Eq:PolFacLemma1} \\
    U \sqrt{Q^{\frac 12} P Q^{\frac 12}} U^* &= \sqrt{P^{\frac 12} Q P^{\frac 12}}, \label{Eq:PolFacLemma2} \\
    U^* &= \operatorname{Pol}\left(Q^{\frac 12} P^{\frac 12}\right), \\
    U &= \operatorname{Pol}\left(P^{-\frac 12} Q^{-\frac 12}\right), \\
    P \cdot (P^{-1} \# Q) &= \sqrt{PQ} = P^{\frac 12} U Q^{\frac 12} \\  
     (P^{-1} \# Q) \cdot P &= \sqrt{QP} = Q^{\frac 12} U^* P^{\frac 12} .
\end{align}
\end{lemma}

\begin{prop} \label{Thm:GBRAlternativeForm}
    Let $P,Q,R \in \mathbb P_d$. Then, the squared generalized Bures distance has the form 
    \begin{equation}
        \operatorname{B}_R(P,Q) = \left\| U^*_P P^{\frac12}  - U^*_Q Q ^{\frac12} \right\|_2^2,
    \end{equation}
    where $U_P := \operatorname{Pol}\left(P^{\frac12}  R ^{\frac12} \right)$, $U_Q := \operatorname{Pol}\left(Q ^{\frac12} R ^{\frac12} \right)$, and $\|A\|_2 := \sqrt{\operatorname{Tr}[A^*A]}$ denotes the Frobenius norm of a matrix $A$. 
\end{prop}

\begin{proof}
    The symmetry part follows directly from the definition. For the nonnegativity part, either observe that the generalized Bures distance is a norm in the tangent space $T_R \mathbb P_d$ or observe that by Proposition~\ref{Thm:GBRAlternativeForm}, we have 
    \begin{equation}
        \operatorname{b}_R(P,Q) = \left\| U_P^* P^{\frac12}  - U_Q^* Q ^{\frac12} \right\|_2 \geq 0, 
    \end{equation}    
    where $U_P := \operatorname{Pol}\left(P^{\frac12}  R  ^{\frac12} \right)$ and $U_Q := \operatorname{Pol}\left(Q ^{\frac12} R  ^{\frac12} \right)$. We now show that equality is achieved if and only if $P = Q$. The direction $P=Q \implies \operatorname{b}_R(P,Q) = 0$ follows trivially. In the other direction, we want to prove
    \begin{equation}
        \operatorname{b}_R(P,Q) = 0 \implies P = Q.
    \end{equation}
    To this end, observe that
    \begin{equation}
        \operatorname{b}_R(P,Q) = \left\| U_P^* P^{\frac12}  - U_Q^* Q^{\frac12} \right\|_2 = 0 \quad \implies  \quad U_P^* P^{\frac12}  =  U_Q^* Q^{\frac12}. 
    \end{equation}
    Right multiply both sides by $R ^{\frac12} $ to obtain
    \begin{equation}
        U_P^* P^{\frac12}  R ^{\frac12} = U_Q^* Q ^{\frac12} R ^{\frac12}. 
    \end{equation}
    By Lemma~\ref{Lem:PolarFactorLemma}, we have that the above equality is the same as
    \begin{equation}
        \sqrt{R ^{\frac12} P R ^{\frac12} } = \sqrt{R ^{\frac12} Q R ^{\frac12} }.
    \end{equation}
    Squaring both sides and right-left multiplying by $ R^{-\frac12} $, we get $P = Q$.
    
    To prove that the generalized Bures distance satisfies the triangle inequality, choose arbitrary $P,Q,R,S \in \mathbb P_d$. We then have,
    \begin{equation}
    \begin{aligned}
        \operatorname{b}_R(P,Q) =  &\left\| U_P^* P^{\frac12}  - U_Q^* Q^{\frac12}   \right\|_2  \\ 
        = &\left\| U_P^* P^{\frac12}  - U_S^* S ^{\frac12} + U_S^* S ^{\frac12} +  U_Q^* Q^{\frac12} \right\|_2 \\ \leq & \left\| U_P^* P^{\frac12}  - U_S^* S ^{\frac12} \right\|_2 + \left\| U_S^* S ^{\frac12} - U_Q^* Q ^{\frac12} \right\|_2  \\ = &\operatorname{b}_R(P,S) + \operatorname{b}_R(S,Q).   
    \end{aligned}
    \end{equation} 
    Here $U_S := \operatorname{Pol}\left(R ^{\frac12} S ^{\frac12} \right)$ and we have used triangle inequality for the Frobenius norm.
\end{proof}

\end{document}
